\documentclass[UTF8,fontset=none,12pt]{ctexart}

\setCJKmainfont{Songti SC}
\setCJKsansfont{PingFang SC}
\setCJKmonofont{Songti SC}

\usepackage[a4paper,margin=2.35cm]{geometry}
\usepackage{amsmath,amssymb,amsthm,mathtools,bm}
\usepackage{graphicx}
\usepackage{float}
\usepackage{booktabs}
\usepackage{array}
\usepackage{enumitem}
\usepackage[hidelinks]{hyperref}

\newcommand{\C}{\mathbb{C}}
\newcommand{\epsmach}{\varepsilon_{\mathrm{machine}}}
\theoremstyle{definition}
\newtheorem{theorem}{定理}
\renewcommand{\thetheorem}{19.\arabic{theorem}}

\title{数值线性代数\quad 第 19 讲：最小二乘算法的稳定性}
\author{Lloyd N. Trefethen、David Bau III\quad 中文阅读版}
\date{}

\begin{document}
\maketitle
\pagestyle{plain}

正如第 11 讲所述，最小二乘问题可以用多种方法求解，包括正规方程、Householder 三角化、Gram--Schmidt 正交化和 SVD.本讲比较这些方法，并说明使用正规方程通常是不稳定的.

\section*{例子}

为了说明各种算法的行为，考虑 $m=100,n=15$ 的数值例子.下面是 MATLAB 设置：

\begin{verbatim}
m = 100; n = 15;
t = (0:m-1)'/(m-1);                 Set t to a discretization of [0, 1].
A = []; for i=1:n,                   Construct Vandermonde matrix.
  A = [A t.^(i-1)]; end
b = exp(sin(4*t));                  Right-hand side.
b = b/2006.787453080206;             Normalization (see text).
\end{verbatim}

这个例子的思想，是在区间 $[0,1]$ 上用 14 次多项式对函数 $\exp(\sin(4t))$ 做最小二乘拟合.首先将 $[0,1]$ 离散化，定义由 100 个等距点组成的向量 $t$.矩阵 $A$ 是 $100\times15$ 的 Vandermonde 矩阵，其列是在这些点处采样的 $1,t,\ldots,t^{14}$；右端 $b$ 是在同一点采样的函数 $\exp(\sin(4t))$.

代码最后一行看起来很奇怪，原因如下.为简单起见，我们只比较不同算法计算出的系数 $x_{15}$.如果没有最后一行，$x_{15}$ 的正确值是
\[
2006.787453080206\ldots,
\]
这个数由扩展精度算术包得到.将 $b$ 除以它之后，问题的解满足 $x_{15}=1$，比较就容易得多.

为了解释观察结果，需要使用式 (18.3)--(18.5) 中的量.可以用 MATLAB 的反斜杠运算符数值求解最小二乘问题，从而得到足够准确的估计：

\begin{verbatim}
x = A\b; y = A*x;                        Solve least squares problem.
kappa = cond(A)
    kappa = 2.2718e+10                   kappa(A)
theta = asin(norm(b-y)/norm(b))
    theta = 3.7461e-06                   theta
eta = norm(A)*norm(x)/norm(y)
    eta = 2.1036e+05                     eta
\end{verbatim}

结果 $\kappa(A)\approx10^{10}$ 表明单项式 $1,t,\ldots,t^{14}$ 构成了一个高度病态的基.结果 $\theta\approx10^{-6}$ 表明 $\exp(\sin(4t))$ 可以被 14 次多项式非常精确地拟合.拟合如此接近，以至于为避免消去误差，我们用
\[
\theta=\sin^{-1}\left(\frac{\|b-y\|}{\|b\|}\right)
\]
而不是式 (18.4) 来计算 $\theta$.至于 $\eta$，它约为 $10^5$，大致位于式 (18.6) 允许的两个极端值 1 和 $\kappa(A)$ 的中间.

将这些数代入定理 18.1，可知在本例中，$b$ 和 $A$ 的扰动对 $y$ 和 $x$ 的条件数近似为
\[
\begin{array}{c|cc}
 & y & x\\
\hline
b&1.0&1.1\times10^5\\
A&2.3\times10^{10}&3.2\times10^{10}
\end{array}
\]

\section*{Householder 三角化}

正如第 11 讲所述，求解最小二乘问题的标准算法，是通过 Householder 三角化进行 QR 分解（算法 11.2）.下面是 MATLAB 实验：

\begin{verbatim}
[Q,R] = qr(A,0);                    Householder triang. of A.
x = R\(Q'*b);                       Solve for x.
x(15)
    ans = 1.00000031528723
\end{verbatim}

如何理解这个结果？由于我们做了归一化，正确答案应为 $x_{15}=1$，所以相对误差约为 $3\times10^{-7}$.计算使用 IEEE 双精度算术，$\epsmach\approx10^{-16}$，这意味着舍入误差被放大了约 $10^9$ 倍.乍看之下这似乎很糟，但上面的表提醒我们，$A$ 的扰动对 $x$ 的条件数约为 $10^{10}$.因此，$x_{15}$ 的不准确完全可以由病态性解释，而不是由不稳定性解释.算法 11.2 看起来是后向稳定的.

上面显式形成了 $Q$，但正如第 10、16 讲强调的，这并非必要.只需存储算法 10.1 每一步确定的向量 $v_k$（式 (10.5)），再用算法 10.2 计算 $Q^*b$.在 MATLAB 中，可以不只对 $A$ 做 QR 分解，而是对 $m\times(n+1)$ 的增广矩阵 $[A\ b]$ 做 QR 分解.在这个过程中，使 $A$ 变为上三角的 $n$ 个 Householder 反射器也会作用于 $b$，把向量 $Q^*b$ 留在第 $n+1$ 列的前 $n$ 个位置.随后应用第 $n+1$ 个反射器，将该列的第 $n+2$ 到第 $m$ 个条目变为零；这不会改变我们关心的前 $n$ 个条目.因此：

\begin{verbatim}
[Q2,R2] = qr([A b],0);              Householder triang. of [A b].
R2 = R2(1:n,1:n);                   Extract R ...
Qb = R2(1:n,n+1);                   ... and Q'*b.
x = R2\Qb;                          Solve for x.
x(15)
    ans = 1.00000031529465
\end{verbatim}

这个答案与前一个几乎相同，说明 QR 分解 $A$ 引入的误差淹没了计算 $Q^*b$ 引入的误差.

还有第三种通过 Householder 三角化求解最小二乘问题的 MATLAB 方法：使用内置反斜杠运算符，正如前面计算 $\kappa(A)$、$\theta$ 和 $\eta$ 时所做的那样：

\begin{verbatim}
x = A\b;                            Solve for x.
x(15)
    ans = 0.99999994311087
\end{verbatim}

这个结果明显不同，而且准确了一个数量级.原因是 MATLAB 的反斜杠运算符使用带列选主元的 QR 分解，其形式为 $AP=QR$，其中 $P$ 是置换矩阵.本书不讨论列选主元.

从范数级稳定性分析的角度看，这三种 QR 分解变体是等价的；可以证明它们全部后向稳定.

\begin{theorem}
设满秩最小二乘问题 (11.2) 在满足公理 (13.5)、(13.7) 的计算机上通过 Householder 三角化（算法 11.2）求解.则该算法后向稳定，即计算得到的解 $\widetilde x$ 满足
\begin{equation}
    \|(A+\delta A)\widetilde x-b\|=\min,
    \qquad \frac{\|\delta A\|}{\|A\|}=O(\epsmach),
    \tag{19.1}
\end{equation}
其中 $\delta A\in\C^{m\times n}$.无论 $Q^*b$ 是通过显式形成 $Q$ 计算，还是由算法 10.2 隐式计算，结论都成立；对任意列选主元的 Householder 三角化也成立.
\end{theorem}

\section*{Gram--Schmidt 正交化}

求解最小二乘问题的另一种方法是改进的 Gram--Schmidt 正交化（算法 8.1）.当 $m\approx n$ 时，它比 Householder 方法需要稍多的运算；但当 $m\gg n$ 时，两种算法的运算量都渐近为 $2mn^2$.

下面的 MATLAB 序列以直接方式实现该算法.函数 `mgs` 是算法 8.1 的一个实现（此处不展示），与第 9 讲实验 2 中的实现相同：

\begin{verbatim}
[Q,R] = mgs(A);                      Gram-Schmidt orthog. of A.
x = R\(Q'*b);                         Solve for x.
x(15)
    ans = 1.02926594532672
\end{verbatim}

这个结果很差.舍入误差被放大了约 $10^{14}$ 倍，远大于问题的条件数.事实上，该算法是不稳定的，原因很容易找出：正如第 9 讲末尾所述，Gram--Schmidt 正交化通常产生列并不准确正交的矩阵 $\widetilde Q$.上面的算法依赖这种正交性，因此会受到相应影响.

可以通过改写算法避免不稳定性.即使 $\widetilde Q$ 的列不准确正交，Gram--Schmidt 迭代仍会给出准确的乘积 $\widetilde Q\widetilde R$.一种办法是为向量 $\widetilde R x$ 建立正规方程
\[
\widetilde R x=(\widetilde Q^*\widetilde Q)^{-1}\widetilde Q^*b,
\]
然后回代求 $x$.只要计算得到的 $\widetilde Q$ 至少是良态的，这种方法就不会遭受下面讨论的、将正规方程应用于任意矩阵时的那些不稳定性.不过它包含不必要的额外工作，实践中不应使用.

稳定 Gram--Schmidt 方法的更好办法，是像前面的第二个 Householder 实验一样使用增广方程组：

\begin{verbatim}
[Q2,R2] = mgs([A b]);                Gram-Schmidt orthog. of [A b].
R2 = R2(1:n,1:n);                    Extract R ...
Qb = R2(1:n,n+1);                    ... and Q'*b.
x = R2\Qb;                           Solve for x.
x(15)
    ans = 1.00000005653399
\end{verbatim}

现在结果与 Householder 三角化一样好.可以证明这总是成立.

\begin{theorem}
满秩最小二乘问题 (11.2) 通过 Gram--Schmidt 正交化求得的解也后向稳定，满足式 (19.1)，前提是像上面的代码片段那样隐式形成 $\widetilde Q^*b$.
\end{theorem}

\section*{正规方程}

求解最小二乘问题的另一种根本不同的方法，是求解正规方程（算法 11.1），通常使用 Cholesky 分解（第 23 讲）.当 $m\gg n$ 时，这种方法比依赖显式正交化的方法快一倍，渐近只需要 $mn^2$ 次浮点运算（式 (11.14)）.在下面的实验中，用一行 MATLAB 反斜杠命令求解：

\begin{verbatim}
x = (A'*A)\(A'*b);                    Form and solve normal equations.
x(15)
    ans = 0.39339069870283
\end{verbatim}

这个结果非常糟糕！这是我们得到的最差结果，甚至没有一位正确数字.显然，使用正规方程求解最小二乘问题是不稳定的.下面解释这一现象；这正是良态性和稳定性相互作用的典型例子.另外，正规方程的使用非常普遍，因此理解其中的风险十分重要.

假设我们有一个求解满秩问题 (11.2) 的后向稳定算法，得到的解 $\widetilde x$ 满足
\[
\|(A+\delta A)\widetilde x-b\|=\min,
\qquad \frac{\|\delta A\|}{\|A\|}=O(\epsmach).
\]
（允许 $b$ 也发生扰动，或者考虑稳定性而非后向稳定性，都不会改变主要结论.）由定理 15.1 和 18.1，有
\begin{equation}
    \frac{\|\widetilde x-x\|}{\|x\|}
    =O\left(\left(\kappa+\frac{\kappa^2\tan\theta}{\eta}\right)\epsmach\right),
    \qquad \kappa=\kappa(A).
    \tag{19.2}
\end{equation}

现在假设 $A$ 病态，即 $\kappa\gg1$，且 $\theta$ 远离 $\pi/2$.根据不同参数的取值，可能出现两种完全不同的情形.如果 $\tan\theta$ 是 1 的量级（即拟合并不特别接近），且 $\eta\ll\kappa$，则式 (19.2) 右端为 $O(\kappa^2\epsmach)$.另一方面，如果 $\tan\theta$ 接近零（拟合非常接近），或者 $\eta$ 接近 $\kappa$，则界为 $O(\kappa\epsmach)$.最小二乘问题的条件数可能处于从 $\kappa$ 到 $\kappa^2$ 的任意位置.

现在考察用正规方程求解 (11.2) 的情形：
\[
(A^*A)x=A^*b.
\]
Cholesky 分解对这个方程组是稳定算法，即它得到的解 $\widetilde x$ 满足
\[
(A^*A+\delta H)\widetilde x=A^*b,
\qquad \frac{\|\delta H\|}{\|A^*A\|}=O(\epsmach)
\]
（定理 23.3）.但是矩阵 $A^*A$ 的条件数是 $\kappa^2$，而不是 $\kappa$.因此正规方程所能期待的最好结果是
\begin{equation}
    \frac{\|\widetilde x-x\|}{\|x\|}=O(\kappa^2\epsmach).
    \tag{19.3}
\end{equation}
正规方程的行为由 $\kappa^2$ 而不是 $\kappa$ 控制.

结论已经很清楚.如果 $\tan\theta$ 为 1 的量级且 $\eta\ll\kappa$，或者 $\kappa$ 为 1 的量级，那么式 (19.2) 和 (19.3) 同阶，正规方程是稳定的.但如果 $\kappa$ 很大，且 $\tan\theta$ 接近零或 $\eta$ 接近 $\kappa$，式 (19.3) 就会远大于式 (19.2)，正规方程是不稳定的.对于包含很接近拟合的病态问题，正规方程通常是不稳定的.在本例中，$\kappa^2\approx10^{20}$，因此 Cholesky 分解得不到任何正确数字并不令人意外.

按我们的定义，只有在所考虑的全部问题上都一致表现良好的算法才是稳定的.因此，可以用下面的结果把刚才的观察正式化.

\begin{theorem}
满秩最小二乘问题 (11.2) 通过正规方程（算法 11.1）得到的解是不稳定的.不过，如果把问题限制在如下某一类中，则可以获得稳定性：$\kappa(A)$ 有统一的上界，或者 $(\tan\theta)/\eta$ 有统一的正下界.
\end{theorem}

\section*{SVD}

第 11 讲还提到过一种最小二乘算法：使用 SVD（算法 11.3）.像大多数基于 SVD 的计算一样，这个方法是稳定的：

\begin{verbatim}
[U,S,V] = svd(A,0);                   Reduced SVD of A.
x = V*(S\(U'*b));                     Solve for x.
x(15)
    ans = 0.99999998230471
\end{verbatim}

事实上，这是所有实验结果中最准确的，比带列选主元的 Householder 三角化（MATLAB 的反斜杠）还准确约 3 倍.可以证明通常形式的定理.

\begin{theorem}
满秩最小二乘问题 (11.2) 通过 SVD（算法 11.3）得到的解后向稳定，满足估计式 (19.1).
\end{theorem}

\section*{秩亏最小二乘问题}

本讲确定了四种用于满秩线性最小二乘问题的后向稳定算法：Householder 三角化、带列选主元的 Householder 三角化、隐式计算 $\widetilde Q^*b$ 的改进 Gram--Schmidt，以及 SVD.从满秩问题 (11.2) 的经典范数级稳定性分析角度看，这些算法之间的差异很小，因此不妨采用最简单、最便宜的不带列选主元的 Householder 三角化.

不过，还有其他类型的最小二乘问题，其中列选主元和 SVD 具有特殊重要性.这类问题满足 $\operatorname{rank}(A)<n$，甚至可能有 $m<n$，因而方程组是不定的.除非增加额外条件（通常要求 $x$ 的范数尽可能小），否则这类问题没有唯一解.另一个困难是，正确解依赖于 $A$ 的秩，而在存在舍入误差时数值确定秩从来都不是小事.

因此，秩亏最小二乘问题不是最小二乘问题的一个简单子类，而是根本不同的问题.既然解的定义已经改变，就没有理由认为对满秩问题稳定的算法在秩亏情形也稳定.事实上，对秩亏问题完全稳定的算法只有基于 SVD 的算法.另一种选择是带列选主元的 Householder 三角化，它对几乎所有问题都稳定.本书不展开细节.

\section*{习题}

\begin{enumerate}[label=\textbf{19.\arabic*}]
    \item 设 $A\in\C^{m\times n}$ 满秩（秩为 $n$），$b\in\C^m$.考虑块 $2\times2$ 方程组
    \begin{equation}
    \begin{bmatrix}I&A\\A^*&0\end{bmatrix}
    \begin{bmatrix}r\\x\end{bmatrix}
    =\begin{bmatrix}b\\0\end{bmatrix},
    \tag{19.4}
    \end{equation}
    其中 $I$ 是 $m\times m$ 单位矩阵.证明这个系统有唯一解 $(r,x)^T$，并证明向量 $r$ 和 $x$ 分别是最小二乘问题 (18.1) 的残差和解.

    \item 下面是 MATLAB 某个内置 m 文件的精简版本：
    \begin{verbatim}
[U,S,V] = svd(A);
S = diag(S);
tol = max(size(A))*S(1)*eps;
r = sum(S > tol);
S = diag(ones(r,1)./S(1:r));
X = V(:,1:r)*S*U(:,1:r)';
    \end{verbatim}
    这个程序计算什么？
\end{enumerate}

\end{document}
